\section*{Bab \ref{ch:b3:lp} --- Ruang $L^p$}

\begin{solution}{exo:b3:lp:1}
(a) Ambil $p = q = 2$ pada \cref{thm:b3:lp:holder}: $\abs{\int f\bar
g\,\dd\mu} \leq \int\abs{fg} \leq \norm f_2\norm g_2$ ---
yakni Cauchy--Schwarz, dengan kesamaannya jika dan hanya jika $\abs f, \abs g$
sebanding dan fasanya sejajar.

(b) Pada ruang peluang, untuk $p \leq q$: terapkan Jensen
(\cref{exo:b3:lp:10}) dengan $\Phi(t) =
\abs t^{q/p}$ yang cembung pada fungsi $\abs f^p$:
maka $\bigl(\int\abs f^p\bigr)^{q/p} \leq \int\abs f^q$, yakni
$\norm f_p \leq \norm f_q$.

(c) Berlaku $\int\abs{fg} = \norm f_1\norm g_\infty$ jika dan hanya jika $\abs{g} =
\norm g_\infty$ hampir di mana-mana pada $\{f \neq 0\}$ (sebab ketaksamaan
$\abs{fg} \leq \abs f\norm g_\infty$ haruslah kesamaan hampir
di mana-mana).
\end{solution}

\begin{solution}{exo:b3:lp:2}
Berlaku $\int_0^1 x^{-p/2}\dd x < \infty$ jika dan hanya jika $p < 2$: sehingga yang pertama berada di
$L^p$ untuk $p \in \intco12$. Lalu $\int_1^\infty x^{-p/2}\dd x <
\infty$ jika dan hanya jika $p > 2$: sehingga yang kedua untuk $p \in \intoo2\infty$ (dan
$p = \infty$: sebab ia terbatas --- sertakanlah). Yang ketiga: untuk $p <
2$, ia terdominasi oleh $x^{-p/2}$: jadi \omterm{def:b3:lebesgue:l1}{terintegralkan}; untuk $p = 2$,
substitusikan $u = -\ln x$: $\int_0^1\frac{\dd x}{x(1 + \abs{\ln
x})^2} = \int_0^\infty\frac{\dd u}{(1 + u)^2} < \infty$; sedangkan untuk
$p > 2$ pangkatnya mendominasi: jadi divergen. Jadi $p \in \intcc12$.
Yang keempat: $\int_\R\frac{\dd x}{(1 + \abs x)^p} < \infty$ jika dan hanya jika $p >
1$; dan ia terbatas, jadi $L^\infty$ pula: sehingga $p \in
\intoc1\infty$. Moralnya: keterintegralan di $0$ menyukai $p$ yang kecil, sedangkan di
$\infty$ $p$ yang besar; dan menggabungkan kedua rintangannya, tak ada keterkandungan
antara ruang $L^p(\R)$.
\end{solution}

\begin{solution}{exo:b3:lp:3}
(a) Berlaku $\norm{f_n}_p^p = \lambda(I_n) \to 0$ (sebab pada aras diadik $k$
panjangnya $2^{-k}$). Tetapi setiap $x$ terletak di satu selang pada
setiap aras diadiknya: sehingga $f_n(x) = 1$ tak berhingga sering dan $= 0$
tak berhingga sering (yakni selang pada aras yang sama yang tak memuat
$x$): jadi tak ada kekonvergenan di titik mana pun.

(b) Barisan $f_{n_k} = \mathbf 1_{\intcc0{2^{-k}}}$ (yakni selang pertama
setiap arasnya) konvergen ke $0$ di setiap $x > 0$: jadi hampir di mana-mana.

(c) Hampir di mana-mana tetapi tak di $L^1$: $n\mathbf 1_{\intoo0{1/n}} \to 0$
hampir di mana-mana, dengan integral $1$. Di $L^1$ tetapi tak di $L^p$ mana pun ($p > 1$): $g_n =
\eu^n\,\mathbf 1_{(0,\ \eu^{-n}/n)}$: sebab $\norm{g_n}_1 = \frac1n
\to 0$, sedangkan $\norm{g_n}_p^p = \eu^{(p-1)n}/n \to \infty$ untuk
setiap $p > 1$.
\end{solution}

\begin{solution}{exo:b3:lp:4}
Batas atasnya: $\norm f_p \leq \mu(X)^{1/p}\norm f_\infty$
(\cref{prop:b3:lp:inclusions}(a) dengan $q = \infty$), dan
$\mu(X)^{1/p} \to 1$. Batas bawahnya: untuk $\varepsilon > 0$, $A =
\{\abs f > \norm f_\infty - \varepsilon\}$ mempunyai $\mu(A) > 0$
(menurut definisi sup esensialnya), dan
\[
\norm f_p \geq \Bigl(\int_A \abs f^p\Bigr)^{1/p}
\geq (\norm f_\infty - \varepsilon)\,\mu(A)^{1/p}
\xrightarrow[p\to\infty]{} \norm f_\infty - \varepsilon .
\]
\end{solution}

\begin{solution}{exo:b3:lp:5}
(a) Dua kali: sebab pengubahan $\norm{\tau_hg - g}_p \leq
C^{1/p}\norm{\tau_hg - g}_\infty$ (dengan pendukung berukuran berhingga)
merosot untuk $p = \infty$ hanya karena \emph{\omterm{ex:b3:lebesgue:gamma}{kepadatan}
$\mathcal C_c$} gagal di sana --- dan itulah jurang yang sebenarnya: sebab langkah (2)
\cref{thm:b3:lp:density} tak beranalogi $L^\infty$.

(b) Jika $f$ mempunyai wakil $g$ yang \omterm{def:b3:topology:continuity}{kontinu} seragam:
maka $\norm{\tau_hf - f}_\infty = \sup_x\abs{g(x - h) - g(x)} \to
0$. Sebaliknya, andaikan $\norm{\tau_hf - f}_\infty \to 0$. Maka
pemulusannya $f_\varepsilon = f * \rho_\varepsilon$ bersifat
\omterm{def:b3:topology:continuity}{kontinu}, dan
\[
\norm{f_\varepsilon - f}_\infty
\leq \sup_{\norm y \leq \varepsilon}\norm{\tau_yf -
f}_\infty \longrightarrow 0
\]
(sebab konvolusinya merupakan rata-rata translasinya). Setiap
$f_\varepsilon$ \omterm{def:b3:topology:continuity}{kontinu} seragam
(sebab $\norm{\tau_hf_\varepsilon - f_\varepsilon}_\infty \leq
\norm{\tau_hf - f}_\infty$, lewat perata-rataannya), dan limit seragam
fungsi \omterm{def:b3:topology:continuity}{kontinu} seragam tetap demikian: sehingga $f$ sama hampir di mana-mana
dengan sebuah fungsi yang \omterm{def:b3:topology:continuity}{kontinu} seragam.
\end{solution}

\begin{solution}{exo:b3:lp:6}
Menurut Hölder, untuk setiap $x$ integrannya $y \mapsto f(x -
y)g(y)$ berada di $L^1$ dengan $\abs{(f*g)(x)} \leq \norm f_p\norm
g_q$: jadi terdefinisi di mana-mana dan terbatas. \omterm{def:b3:topology:continuity}{Kekontinuan} seragamnya (bila $p <
\infty$):
\[
\abs{(f*g)(x + h) - (f*g)(x)}
= \Bigl|\int\bigl(\tau_{-h}f - f\bigr)(x - y)\,g(y)\dd y\Bigr|
\leq \norm{\tau_{-h}f - f}_p\,\norm g_q
\xrightarrow[h\to0]{} 0,
\]
secara seragam terhadap $x$ (\cref{thm:b3:lp:density}(3)). Jika $p =
\infty$, maka $q = 1$: tulislah $f * g = g * f$ lalu jalankan batas yang
sama dengan translasinya bekerja pada $g \in L^1$.
\end{solution}

\begin{solution}{exo:b3:lp:7}
Tetapkan $\chi \in \mathcal C_c^\infty(\intoo ab)$ dengan $\int\chi =
1$, lalu tetapkan $c = \int f\chi$. Misalkan $\varphi \in \mathcal
C_c^\infty$ sembarang dan $\psi = \varphi -
\bigl(\int\varphi\bigr)\chi$: maka $\int\psi = 0$, sehingga $\Phi(x)
= \int_a^x\psi$ mendefinisikan $\Phi \in \mathcal C_c^\infty(\intoo
ab)$ (sebab ia lenyap di dekat kedua ujungnya: di dekat $a$ secara trivial, dan di dekat $b$
karena integral totalnya $0$) dengan $\Phi' = \psi$. Lalu
hipotesisnya memberikan $\int f\psi = \int f\Phi' = 0$, sehingga
\[
\int f\varphi = \Bigl(\int\varphi\Bigr)\int f\chi = \int
c\,\varphi
\quad\text{untuk setiap } \varphi:
\qquad \int(f - c)\varphi = 0 .
\]
Menurut \cref{cor:b3:lp:fundlemma} (yang dilokalkan pada $\intoo ab$), $f =
c$ hampir di mana-mana.
\end{solution}

\begin{solution}{exo:b3:lp:8}
Misalkan $3\delta < d(K, \R^d\setminus U)$ (yang positif:
\cref{exo:b3:topology:6}(b); dan jika $U = \R^d$ maka sembarang $\delta$
berhasil), $K_\delta = \{x : d(x, K) \leq \delta\}$, dan
\[
\varphi = \mathbf 1_{K_\delta} * \rho_{\delta/2} .
\]
Maka $\varphi \in \mathcal C^\infty$
(\cref{thm:b3:lp:regularization}(1); sebab indikatornya $L^1$),
$0 \leq \varphi \leq 1$ (sebab $\int\rho = 1$), $\varphi = 1$ pada $K$
(sebab untuk $x \in K$, $\bar B(x, \delta/2) \subseteq K_\delta$, sehingga
konvolusinya mengintegralkan $\rho$ sepenuhnya), dan
$\operatorname{supp}\varphi \subseteq K_{3\delta/2} \subseteq
U$: jadi berpendukung \omterm{def:b3:topology:compact}{kompak} (sebab $K_\delta$ terbatas). Untuk partisi
satuannya: diberikan $K \subseteq U_1\cup\dots\cup U_m$, pilihlah (lewat
kekompakannya) \omterm{def:b3:topology:compact}{kompak} $K_i \subseteq U_i$ dengan $K \subseteq
\bigcup \mathring K_i$, ambil $\varphi_i$ seperti di atas untuk $(K_i,
U_i)$, lalu tetapkan $\psi_i = \varphi_i\prod_{j <i}(1 -
\varphi_j)$: maka setiap $\psi_i \in \mathcal C_c^\infty(U_i)$, dan
$\sum_i\psi_i = 1 - \prod_i(1 - \varphi_i) = 1$ pada $K$.
\end{solution}

\begin{solution}{exo:b3:lp:9}
(a) Eksponen interpolasinya untuk $(p_0, q_0) = (1, \infty)$
di $r = p$ adalah $\alpha = \frac1p$: sehingga
\cref{prop:b3:lp:inclusions}(c) memberikan $\norm f_p \leq \norm
f_1^{1/p}\norm f_\infty^{1 - 1/p}$.

(b) Dengan mengambil logaritma pada
\cref{prop:b3:lp:inclusions}(c): $\ln\norm f_r \leq
\alpha\ln\norm f_p + (1-\alpha)\ln\norm f_q$ dengan $\frac1r$
berupa kombinasi cembung yang sama atas $\frac1p, \frac1q$: sehingga $\frac1p
\mapsto \ln\norm f_p$ cembung. Contoh yang keanggotaan $L^p$-nya
tepat pada $\intoo{p_0}{p_1}$:
\[
f(x) = x^{-1/p_1}\,\mathbf 1_{\intoo01}(x) +
x^{-1/p_0}\,\mathbf 1_{\intco1\infty}(x):
\]
sebab suku pertamanya $L^p$ jika dan hanya jika $p < p_1$, dan yang kedua jika dan hanya jika $p >
p_0$.
\end{solution}

\begin{solution}{exo:b3:lp:10}
Misalkan $m = \int f\,\dd\mu \in \R$. Kecembungannya menyediakan garis
penopang di $m$: yakni ada $s$ dengan $\Phi(t) \geq \Phi(m) + s(t -
m)$ untuk setiap $t$ (ambillah $s$ di antara turunan sepihaknya,
yang ada bagi fungsi cembung). Lalu substitusikan $t = f(x)$ dan
integralkan terhadap peluang $\mu$:
\[
\int\Phi\circ f\,\dd\mu \geq \Phi(m) + s\Bigl(\int f - m\Bigr)
= \Phi\Bigl(\int f\,\dd\mu\Bigr)
\]
(\omterm{def:b3:lebesgue:measurable}{keterukurannya}: sebab $\Phi$ \omterm{def:b3:topology:continuity}{kontinu}; sedangkan keterintegralan bagian
negatif $\Phi\circ f$ dijamin oleh garis
penopangnya). Untuk rata-rata hitung--ukur: pada himpunan berhingga berbobot $w_i$, ambillah
$\Phi = \exp$ dan $f = \sum(\ln a_i)\mathbf 1_i$:
maka $\exp\bigl(\sum w_i\ln a_i\bigr) \leq \sum w_ia_i$, yakni
$\prod a_i^{w_i} \leq \sum w_ia_i$. Sedangkan kemonotonan normanya adalah
\cref{exo:b3:lp:1}(b).
\end{solution}

\begin{solution}{exo:b3:lp:11}
(a) Andaikan dahulu $p, q, r < \infty$ dan $f, g \geq 0$
(gantilah dengan nilai mutlaknya). Ketiga eksponennya $r$,
$\alpha = \frac{pr}{r-p}$, dan $\beta = \frac{qr}{r-q}$
memenuhi $\frac1r + \frac1\alpha + \frac1\beta = \frac1r +
\frac1p - \frac1r + \frac1q - \frac1r = 1$ (yakni hubungan
penskalaannya). Lalu pecahlah, untuk $x$ yang tetap,
\[
f(y)g(x-y) = \bigl[f(y)^pg(x-y)^q\bigr]^{1/r}\cdot
f(y)^{1 - p/r}\cdot g(x-y)^{1 - q/r},
\]
dan Hölder dengan ketiga eksponennya memberikan
\[
(f*g)(x) \leq \Bigl(\int f^pg(x-\cdot)^q\Bigr)^{1/r}
\norm f_p^{\,p/\alpha\cdot\alpha/p}\cdots
\]
lebih tepatnya: faktor keduanya adalah
$\bigl(\int f^{(1-p/r)\alpha}\bigr)^{1/\alpha} =
\norm f_p^{p(1/p - 1/r)}$ sebab $(1 - \frac pr)\alpha = p$,
dan demikian pula yang ketiga adalah $\norm g_q^{q(1/q - 1/r)}$. Lalu pangkatkan
ke $r$ dan integralkan terhadap $x$ (dengan Tonelli pada
faktor pertamanya):
\[
\norm{f*g}_r^r \leq \norm f_p^p\,\norm g_q^q\cdot
\norm f_p^{\,rp(1/p - 1/r)}\,\norm g_q^{\,rq(1/q - 1/r)}
= \norm f_p^r\,\norm g_q^r .
\]
Sedangkan kasus ujungnya ($r = \infty$ atau sebuah eksponen yang sama dengan
batasnya) berupa Hölder biasa atau taksiran langsung.

(b) $r = \infty$ memaksa $q = p'$: sehingga $\abs{f*g(x)} \leq \norm
f_p\norm g_{p'}$ --- yakni Hölder setelah translasi dan pencerminan.
Lalu $q = 1$ memberikan $r = p$: $\norm{f*g}_p \leq \norm
g_1\norm f_p$, yakni kuda beban pemulusannya
(yaitu mesin \cref{thm:b3:lp:regularization}). Sedangkan $p = q = 1$
memberikan $r = 1$: yakni aljabar konvolusinya
(\cref{thm:b3:product:convolution}).

(c) Gantilah $f, g$ dengan $f_\lambda = f(\lambda\cdot)$ dan
$g_\lambda = g(\lambda\cdot)$: maka $f_\lambda*g_\lambda =
\lambda^{-d}(f*g)(\lambda\cdot)$, dan dengan membandingkan normanya,
\[
\text{ruas kiri} \sim \lambda^{-d - d/r}, \qquad
\text{ruas kanan} \sim \lambda^{-d/p - d/q} :
\]
sehingga ketaksamaan yang sahih bagi setiap $f, g$ memaksa kedua eksponen
penskalaannya bersesuaian, yakni $1 + \frac1r = \frac1p + \frac1q$
secara persis. Sedangkan gabungan lain mana pun mati saat $\lambda \to 0$ atau
$\infty$: jadi hubungan Young bukanlah kemudahan melainkan sebuah
hukum penskalaan.
\end{solution}

\begin{solution}{exo:b3:lp:12}
(a) Normalkan $\norm f_p = \norm g_q = 1$. Bukti
Hölder mengintegralkan ketaksamaan Young $\abs{fg} \leq
\frac{\abs f^p}p + \frac{\abs g^q}q$; dan kesamaan
integralnya memaksa kesamaan hampir di mana-mana pada Young, yang (lewat kecembungan
sejati $\exp$; dengan kesamaannya jika dan hanya jika $a^p = b^q$) berarti $\abs
f^p = \abs g^q$ hampir di mana-mana. Setelah penormalannya dibatalkan:
$\abs f^p\norm g_q^q = \abs g^q\norm f_p^p$ hampir di mana-mana ---
yakni kesebandingannya.

(b) Minkowski berupa dua Hölder yang diterapkan pada $\abs{f + g}^{p-1}
\abs f$ dan $\abs{f+g}^{p-1}\abs g$; sehingga kesamaannya memaksa
kesebandingan (a) pada keduanya: yakni $\abs f^p$ dan $\abs g^p$ masing-masing
sebanding dengan $\abs{f+g}^{(p-1)q} = \abs{f+g}^p$, sehingga
$\abs g = t\abs f$ hampir di mana-mana untuk sebuah konstanta $t \geq 0$; dan
ketaksamaan segitiga titik demi titik awalnya $\abs{f + g} \leq
\abs f + \abs g$ pun haruslah kesamaan hampir di mana-mana, yang untuk
nilai kompleks berarti $f$ dan $g$ berargumen sama hampir
di mana-mana di tempat keduanya tak nol. Digabungkan: $g = tf$ hampir di mana-mana
dengan $t > 0$ (sebab keduanya tak nol).

(c) Untuk $p = 1$: kesamaan pada $\int\abs{f + g} = \int\abs f +
\int\abs g$ berlaku setiap kali $f, g$ berpola tanda yang sama
(yakni berargumen sama hampir di mana-mana) --- tanpa perlu kesebandingan:
sebab $f = \mathbf 1_{\intcc01}$ dan $g = \mathbf 1_{\intcc02}$
berhasil. Sedangkan untuk $p = \infty$: $\norm{f+g}_\infty = \norm f_\infty +
\norm g_\infty$ segera setelah kedua fungsinya memuncak
secara serasi di sebuah titik bersama (atau sepanjang sebuah barisan bersama):
sehingga $f = g\,$ di dekat satu titik sudah cukup, apa pun perilakunya
di tempat lain. Jadi kecembungan sejati bola $L^p$ untuk $1 <
p < \infty$ --- dan kegagalannya di titik ujungnya --- tepat
merupakan hal yang disaksikan kasus kesamaan ini.
\end{solution}

\begin{solution}{pb:b3:lp:1}
\textbf{1.} Fungsi $f$ berpendukung di suatu $[\alpha, \beta]
\subseteq \intoo0\infty$, sehingga $F = 0$ pada $[0, \alpha]$ dan $F
\equiv F(\beta)$ pada $[\beta, \infty)$: jadi $Hf$ lenyap di dekat $0$
dan bernilai $O(1/x)$ di tak hingga; sedangkan $\int_1^\infty x^{-p}\dd x <
\infty$ untuk $p > 1$, dan $Hf$ \omterm{def:b3:topology:continuity}{kontinu}: sehingga $Hf \in L^p$.

\textbf{2.} Berlaku $\bigl(x^{1-p}F(x)^p\bigr)' = (1-p)x^{-p}F^p +
p\,x^{1-p}F^{p-1}f$. Kedua nilai batasnya lenyap: di $0$
karena $F = 0$ di dekat $0$; dan di $\infty$ karena $x^{1-p}F^p \leq
F(\beta)^p x^{1-p} \to 0$ (sebab $p > 1$). Lalu dengan mengintegralkan identitasnya
atas $\intoo0\infty$:
\[
0 = (1 - p)\int_0^\infty\Bigl(\frac Fx\Bigr)^p\dd x
+ p\int_0^\infty\Bigl(\frac Fx\Bigr)^{p-1}f(x)\,\dd x,
\]
yang merupakan hubungan yang ditampilkan itu.

\textbf{3.} Hölder dengan eksponen $q = \frac p{p-1}$ dan $p$:
\[
\int\Bigl(\frac Fx\Bigr)^{p-1}f
\leq \Bigl(\int\Bigl(\frac
Fx\Bigr)^{p}\Bigr)^{1 - 1/p}\,\norm f_p ,
\]
sehingga $\norm{Hf}_p^p \leq \frac p{p-1}\norm{Hf}_p^{p-1}\norm
f_p$; lalu bagilah (yang berhingga menurut pertanyaan 1, dan jika $0$ maka tak ada
yang perlu dibuktikan).

\textbf{4.} Misalkan $f \in L^p$ dengan $f \geq 0$. Pilihlah $\varphi_k \in
\mathcal C_c(\intoo0\infty)$ dengan $\varphi_k \to f$ di $L^p$
(\cref{thm:b3:lp:density}(2), yang diiriskan dengan setengah garis
terbukanya --- hampirilah $f\mathbf 1_{[1/k, k]}$ lalu
diagonalkan), lalu gantilah $\varphi_k$ dengan $\abs{\varphi_k}$
(yang tetap \omterm{def:b3:topology:continuity}{kontinu}, dan lebih dekat ke $f \geq 0$). Untuk setiap
$x > 0$ yang tetap, Hölder pada $\intoo0x$ memberikan
\[
\abs{H\varphi_k(x) - Hf(x)} \leq
\frac1x\,x^{1 - 1/p}\,\norm{\varphi_k - f}_p \to 0 :
\]
sehingga $H\varphi_k \to Hf$ secara titik demi titik. Lalu Fatou dan pertanyaan 3:
\[
\int (Hf)^p \leq \liminf_k\int(H\varphi_k)^p
\leq \Bigl(\frac p{p-1}\Bigr)^p\liminf_k\norm{\varphi_k}_p^p
= \Bigl(\frac p{p-1}\Bigr)^p\norm f_p^p .
\]
Sedangkan untuk $f$ yang bertanda atau kompleks: $\abs{Hf} \leq H\abs f$ secara titik demi titik,
dan kasus tak negatifnya berlaku bagi $\abs f$.

\textbf{5.} Berlaku $\norm{f_A}_p^p = \int_1^A t^{-1}\dd t = \ln A$.
Untuk $1 \leq x \leq A$:
\[
(Hf_A)(x) = \frac1x\int_1^x t^{-1/p}\dd t
= \frac{x^{1 - 1/p} - 1}{(1 - \tfrac1p)\,x}
= \frac p{p-1}\,x^{-1/p}\bigl(1 - x^{-(1 - 1/p)}\bigr).
\]

\textbf{6.} Tetapkan $\varepsilon > 0$ dan $X_0$ dengan $(1 -
x^{-(1-1/p)})^p \geq 1 - \varepsilon$ untuk $x \geq X_0$. Maka
\[
\norm{Hf_A}_p^p \geq \int_{X_0}^A\Bigl(\frac
p{p-1}\Bigr)^p\frac{1 - \varepsilon}{x}\,\dd x
= \Bigl(\frac p{p-1}\Bigr)^p(1 - \varepsilon)\,(\ln A - \ln
X_0),
\]
sehingga $\dfrac{\norm{Hf_A}_p^p}{\norm{f_A}_p^p} \geq \bigl(\frac
p{p-1}\bigr)^p(1 - \varepsilon)\bigl(1 - \frac{\ln X_0}{\ln
A}\bigr) \to \bigl(\frac p{p-1}\bigr)^p(1 - \varepsilon)$ saat
$A \to \infty$: jadi konstantanya tak dapat diperbaiki.

\textbf{7.} Kesamaan pada pertanyaan 3 memaksa kesamaan pada
Hölder: yakni $f^p$ sebanding dengan $\bigl(\frac Fx\bigr)^{(p-1)q}
= \bigl(\frac Fx\bigr)^p$ hampir di mana-mana, yakni $f = \gamma\,\frac Fx$
hampir di mana-mana untuk suatu $\gamma \geq 0$. Karena $F(x) = \int_0^xf$
\omterm{def:b3:topology:continuity}{kontinu} mutlak dengan $F' = f$ hampir di mana-mana, maka $F$ memenuhi $F' =
\gamma F/x$: sehingga pada sembarang selang yang $F > 0$-nya, $(\ln F)' =
\gamma/x$, jadi $F = c\,x^{\gamma}$ dan $f = c\gamma
x^{\gamma - 1}$ di sana. Padahal tak ada pangkat tak nol $x^{\gamma-1}$ yang
termasuk $L^p(\intoo0\infty)$ (sebab $p(\gamma - 1) < -1$ dituntut
di $\infty$ dan $> -1$ di $0$: tak serasi), dan $F$ tak dapat
lenyap secara identik kecuali $f = 0$. Jadi kesamaannya menuntut $f =
0$.

\textbf{8.} Diberikan $(a_n)$ tak naik dengan $\geq 0$, definisikan
fungsi tangga $g(t) = a_{\lceil t\rceil}$ pada
$\intoo0{+\infty}$: yang tak naik, dengan $\int_0^\infty g^p =
\sum_na_n^p$ dan $\int_0^ng = a_1 + \dots + a_n$, sehingga
$(Hg)(n) = \frac{a_1 + \dots + a_n}n$. Rata-rata sebuah fungsi
tak naik tetap tak naik, sehingga $Hg$ demikian, dan
\[
\sum_{n\geq1}\Bigl(\frac{a_1{+}\dots{+}a_n}n\Bigr)^p
= \sum_{n\geq1}(Hg)(n)^p
\leq \sum_{n\geq1}\int_{n-1}^n (Hg)(t)^p\,\dd t
= \norm{Hg}_p^p
\leq \Bigl(\frac p{p-1}\Bigr)^p\sum_n a_n^p
\]
menurut Bagian I. Sedangkan untuk barisan tak negatif yang umum, misalkan $(a_n^*)$
penyusunan ulangnya yang tak naik (yang mungkin bila $a_n \to 0$,
yang boleh kita andaikan --- sebab jika tidak kedua ruasnya tak berhingga):
maka ruas kanannya tak berubah, dan setiap jumlah parsial $a_1 + \dots
+ a_n$ paling banyak $a_1^* + \dots + a_n^*$ (yakni $n$ suku terbesarnya): sehingga
ruas kirinya hanya bertambah. Karena itu ketaksamaannya berlaku bagi setiap
$(a_n)$.

\textbf{9.} Jika $(a_n) \in \ell^p$, maka barisan rata-rata Cesàro
berada di $\ell^p$ dengan norma $\leq \frac p{p-1}\norm
a_p$. Contohnya ($p = 2$): $a_n = \frac1{\sqrt n\,\ln n}$ (untuk $n
\geq 2$): sebab $\sum a_n^2 = \sum\frac1{n\ln^2n} < \infty$, padahal
$\sum a_n = \infty$ (lewat uji integralnya); dan Hardy tetap menjamin
$\sum_n\bigl(\frac{a_1 + \dots + a_n}n\bigr)^2 < \infty$.

\textbf{10.} Untuk $f = \mathbf 1_{\intcc01}$: $Hf(x) = 1$ pada
$\intoc01$ dan $= \frac1x$ untuk $x \geq 1$: sehingga $\int_0^\infty Hf =
1 + \int_1^\infty\frac{\dd x}x = \infty$, padahal $\norm f_1 =
1$. Buktinya runtuh di dua tempat: sebab konstanta
$\frac p{p-1}$ meledak saat $p \to 1$, dan suku batasnya
$x^{1-p}F^p$ tak lagi lenyap di tak hingga bila $p = 1$.
Jadi ketaksamaan Hardy sungguh gejala $p > 1$.

\textbf{11.} Pilihlah selang terpanjang $B_{i_1}$; buanglah
setiap selang yang memotongnya; pilih penyintas terpanjang
$B_{i_2}$; lalu iterasikan (sebab selangnya berhingga banyak). Yang terpilih
saling lepas menurut konstruksinya, dan setiap $B$ yang terbuang memotong
sebuah selang terpilih yang setidaknya sepanjang itu: dan selang yang memotong selang
yang lebih panjang atau sama termuat di lipat tiganya, yakni $B \subseteq
3B_{i_l}$.

\textbf{12.} Himpunan $\{Mf > t\}$ terbuka: sebab setiap rata-rata $x \mapsto
\frac1{2r}\int_{x-r}^{x+r}\abs f$ \omterm{def:b3:topology:continuity}{kontinu} (lewat kekonvergenan
terdominasi terhadap $x$), dan supremum fungsi \omterm{def:b3:topology:continuity}{kontinu}
bersifat semikontinu bawah. Lalu setiap $x$ di dalamnya memiliki $B_x =
\intoo{x-r_x}{x+r_x}$ dengan $\int_{B_x}\abs f >
t\,\lambda(B_x)$. Untuk \omterm{def:b3:topology:compact}{kompak} $K \subseteq \{Mf > t\}$:
berhingga banyak $B_x$ menyelimuti $K$, lalu Vitali (pertanyaan 11) memetik
$B_1', \dots, B_k'$ yang saling lepas dengan $K \subseteq
\bigcup_l3B_l'$, sehingga
\[
\lambda(K) \leq 3\sum_l\lambda(B_l') <
\frac3t\sum_l\int_{B_l'}\abs f \leq \frac3t\,\norm f_1
\]
lewat kelepasannya; lalu keteraturan dalamnya
(\cref{thm:b3:measure:regularity}) menuntaskannya.

\textbf{13.} Untuk $x > 1$: dengan $r \in \intcc{x-1}x$ rata-ratanya
adalah $\frac{r - x + 1}{2r}$, yang naik terhadap $r$; sedangkan untuk $r
\geq x$ ia $\frac1{2r}$, yang turun: jadi supremumnya
$\frac1{2x}$, yang tercapai di $r = x$. Sehingga $M\mathbf 1_{\intcc01}
\notin L^1$. Secara umum, jika $\int_I\abs f = c > 0$ pada sebuah selang
terbatas $I \subseteq \intcc{-C}C$, maka $Mf(x) \geq
\frac{c}{2(\abs x + C)}$ untuk setiap $x$: jadi tak pernah \omterm{def:b3:lebesgue:l1}{terintegralkan}
kecuali $f = 0$ hampir di mana-mana.

\textbf{14.} Dengan $f_t = f\,\mathbf 1_{\abs f > t/2}$:
$M(f - f_t) \leq \frac t2$, sehingga $\{Mf > t\} \subseteq \{Mf_t
> \frac t2\}$ dan pertanyaan 12 memberikan $\lambda(Mf > t) \leq
\frac6t\int_{\abs f > t/2}\abs f$. Lalu kue berlapis
(\cref{prop:b3:product:layercake}) dan Tonelli:
\[
\norm{Mf}_p^p = p\int_0^\infty t^{p-1}\lambda(Mf > t)\dd t
\leq 6p\int\abs f\int_0^{2\abs f}t^{p-2}\,\dd t\,\dd\lambda
= \frac{6p\,2^{p-1}}{p-1}\,\norm f_p^p .
\]

\textbf{15.} Tulis $A_rf(x) = \frac1{2r}\int_{x-r}^{x+r}f$.
Untuk $g$ yang \omterm{def:b3:topology:continuity}{kontinu}: $A_rg(x) \to g(x)$ di mana-mana. Diberikan
$\varepsilon > 0$, pecahlah $f = g + h$ dengan $g$ \omterm{def:b3:topology:continuity}{kontinu} berpendukung
\omterm{def:b3:topology:compact}{kompak} dan $\norm h_1 < \varepsilon$
(\cref{thm:b3:lp:density}); maka
\[
\limsup_{r\to0}\,\abs{A_rf(x) - f(x)} \leq Mh(x) +
\abs{h(x)},
\]
sehingga $\Omega_\delta = \{\limsup_r\abs{A_rf - f} > \delta\}
\subseteq \{Mh > \tfrac\delta2\} \cup \{\abs h >
\tfrac\delta2\}$ berukuran $\leq \frac{6\varepsilon}\delta
+ \frac{2\varepsilon}\delta$ (lewat pertanyaan 12; dan Markov).
Karena $\varepsilon$ sembarang: $\lambda(\Omega_\delta) = 0$; lalu gabungkan
atas $\delta = \frac1k$: sehingga $A_rf \to f$ hampir di mana-mana.

\textbf{16.} (a) Untuk setiap $q \in \Q$, pertanyaan 15 yang diterapkan pada
$\abs{f - q}$ memberikan $A_r\abs{f - q}(x) \to \abs{f(x) - q}$
hampir di mana-mana; lalu pada irisan himpunan berukuran penuh itu, pilihlah
$q$ dengan $\abs{f(x) - q} < \eta$:
maka $\limsup_rA_r\abs{f - f(x)}(x) \leq 2\eta$ untuk setiap $\eta$: sehingga
hampir setiap $x$ merupakan titik Lebesgue. (b) Pada sebuah titik
Lebesgue,
\[
\Bigl|\frac{F(x + h) - F(x)}h - f(x)\Bigr|
= \Bigl|\frac1h\int_x^{x+h}\bigl(f - f(x)\bigr)\Bigr|
\leq 2\,A_{\abs h}\abs{f - f(x)}(x) \to 0 :
\]
sehingga $F' = f$ hampir di mana-mana --- jadi antiturunan fungsi $L^1$ sungguh
dapat diturunkan kembali; sedangkan tangganya (\cref{pb:b3:measure:1})
merupakan contoh penyangkal bagi arah \emph{sebaliknya} saja.

\textbf{17.} Terapkan pertanyaan 15 pada $\mathbf 1_{A\cap[-n,n]}$
lalu biarkan $n$ tumbuh: maka untuk hampir setiap $x \in A$ \omterm{ex:b3:lebesgue:gamma}{kepadatannya}
$\frac{\lambda(A\cap\intcc{x-r}{x+r})}{2r} \to 1$. Untuk Steinhaus:
di sekitar sebuah titik \omterm{ex:b3:lebesgue:gamma}{kepadatan} ambillah $r$ yang berkepadatan $> \frac34$;
lalu untuk $\abs t < \frac r2$, $A$ dan $A + t$ masing-masing mengisi lebih dari
$\frac32r$ sebuah selang berpanjang $\leq \frac52r$, sehingga
keduanya berpotongan: jadi $\intoo{-r/2}{r/2} \subseteq A - A$.

\textbf{18.} Misalkan $G(x) = x^{\alpha+1-p}F(x)^p$: maka $G$ lenyap
di $0$ (sebab $F$ lenyap di dekat $0$) dan di $\infty$ (sebab $F$ terbatas dan
$\alpha + 1 - p < 0$), sehingga $\int_0^\infty G' = 0$ dengan
\[
G' = (\alpha + 1 - p)\,x^{\alpha-p}F^p +
p\,x^{\alpha+1-p}F^{p-1}f
= (\alpha + 1 - p)\Bigl(\frac Fx\Bigr)^px^\alpha +
p\Bigl(\frac Fx\Bigr)^{p-1}f\,x^\alpha .
\]
Karena itu $\int(F/x)^px^\alpha = \frac{p}{p-1-\alpha}
\int(F/x)^{p-1}f\,x^\alpha$; lalu Hölder untuk \omterm{def:b3:measure:measure}{ukuran}
$x^\alpha\dd x$ (dengan eksponen $\frac p{p-1}$ dan $p$) menyelesaikannya
seperti pada pertanyaan 3. Untuk perbatasan $\alpha = p - 1$: dengan $f(t) =
\frac1t\mathbf 1_{\intcc1A}$, ruas kanannya adalah $\ln A$
sedangkan ruas kirinya memuat $\int_1^A\frac{(\ln x)^p}x\dd x =
\frac{(\ln A)^{p+1}}{p+1}$: sehingga tak ada konstanta yang selamat saat $A \to
\infty$.

\textbf{19.} Tonelli pada $\{0 < t < x\}$:
\[
\langle Hf, g\rangle =
\int_0^\infty\frac{g(x)}x\int_0^xf(t)\,\dd t\,\dd x
= \int_0^\infty f(t)\int_t^\infty\frac{g(x)}x\,\dd x\,\dd t
= \langle f, H^*g\rangle .
\]
Untuk dualitasnya: $\norm{H^*f}_p = \sup\{\langle f, Hg\rangle : g
\geq 0, \norm g_q \leq 1\} \leq \norm f_p\cdot
\sup\norm{Hg}_q \leq \frac q{q-1}\norm f_p = p\,\norm f_p$,
sebab Hardy dipanggil di $L^q$, yang konstantanya $\frac q{q-1}$
sama dengan $p$.

\textbf{20.} Pisahkan sepanjang diagonalnya (yang nol). Pada $\{y \leq
x\}$:
\[
\iint_{y\leq x}\frac{f(x)g(y)}{x}\,\dd y\,\dd x
= \int f(x)\,(Hg)(x)\,\dd x \leq \norm f_p\,\norm{Hg}_q
\leq p\,\norm f_p\norm g_q,
\]
sebab Hardy di $L^q$ membawa konstanta $\frac q{q-1} =
p$. Setangkup dengan itu, $\iint_{y>x} = \int g\,(Hf) \leq
\norm g_q\norm{Hf}_p \leq q\,\norm f_p\norm g_q$ (sebab $\frac
p{p-1} = q$). Totalnya: $(p + q)\,\norm f_p\norm g_q$.

\textbf{21.} Untuk $a_n = n^{-1/p}$ dengan $n \leq N$: ruas
kanannya adalah $\bigl(\frac p{p-1}\bigr)^p\sum_{n\leq N}\frac1n =
\bigl(\frac p{p-1}\bigr)^p\ln N + O(1)$. Sedangkan di kirinya, untuk $n
\leq N$: $a_1 + \dots + a_n \geq \int_1^{n+1}t^{-1/p}\dd t =
\frac{p}{p-1}\bigl((n+1)^{1-1/p} - 1\bigr)$, sehingga rata-rata
Cesàro ke-$n$ bernilai $\geq \frac p{p-1}n^{-1/p}(1 -
o(1))$ secara seragam untuk $n$ pada sembarang rentang $n \geq n_0(\eta)$;
lalu dengan memangkatkannya ke $p$ dan menjumlahkannya, ruas kirinya $\geq
\bigl(\frac p{p-1}\bigr)^p(1 - \eta)\ln N + O_\eta(1)$.
Lalu dengan membaginya dan membiarkan $N \to \infty$, lalu $\eta \to 0$: tak ada
konstanta yang lebih kecil daripada $\bigl(\frac p{p-1}\bigr)^p$ yang dapat bekerja.

\textbf{22.} Keterbatasan $H$ pada $L^p$: rata-rata kumulatifnya tak
menggelembungkan norma-$p$ (dengan konstanta $\frac p{p-1}$); untuk Cesàro pada
$\ell^p$: hal yang sama, terdiskretkan; sedangkan keterbatasan $M$ pada $L^p$: bahkan
rata-rata lokal terbaiknya tetap terkendali (dengan konstanta
$O(\frac1{p-1})$). Ketiganya gagal di $p = 1$, dan karena satu
alasan: sebab merata-ratakan satu satuan massa yang terpusat menghasilkan
ekor $\frac1x$ (lihat pertanyaan 10 dan 13), dan $\frac1x$ termasuk
setiap $L^p$ di dekat tak hingga kecuali $L^1$. Jadi pemulusan menyebarkan
massanya tepat sampai perbatasan harmonik keterintegralannya.

\textbf{23.} Tetapkan $b_n = a_n^{1/p}$, sehingga $\sum b_n^p = \sum
a_n$. Lalu pertanyaan 8 memberikan
\[
\sum_{n\geq1}\Bigl(\frac{b_1 + \dots + b_n}n\Bigr)^{p}
\leq \Bigl(\frac p{p-1}\Bigr)^{p}\sum_{n\geq1}a_n,
\]
dan rata-rata hitung--ukur membatasi setiap sukunya dari bawah:
\[
\Bigl(\frac{b_1 + \dots + b_n}n\Bigr)^{p}
\geq \bigl(b_1\cdots b_n\bigr)^{p/n}
= \bigl(a_1\cdots a_n\bigr)^{1/n}.
\]
Karena itu $\sum(a_1\cdots a_n)^{1/n} \leq \bigl(\frac
p{p-1}\bigr)^p\sum a_n$ untuk \emph{setiap} $p > 1$. Dengan $m =
p - 1$,
\[
\Bigl(\frac p{p-1}\Bigr)^{p} =
\Bigl(1 + \frac1m\Bigr)^{m+1},
\]
yang turun menuju $\eu$ saat $m \to \infty$ (yakni barisan atas
monoton klasik bagi $\eu$). Lalu mengambil infimumnya atas
$p$ memberikan ketaksamaan Carleman berkonstanta $\eu$.

\textbf{24.} Untuk $a_n = \frac1n$, $(a_1\cdots a_n)^{1/n} =
(n!)^{-1/n}$. Pengapitan \cref{pb:b3:product:1} memberikan
$\sqrt{2\pi n}\,(n/\eu)^n \leq n! \leq \sqrt{2\pi
n}\,(n/\eu)^n\eu^{1/(12n)}$, sehingga
\[
(n!)^{1/n} = \frac n\eu\,(2\pi n)^{1/(2n)}
\eu^{O(1/n^2)}
= \frac n\eu\Bigl(1 +
O\Bigl(\frac{\ln n}{n}\Bigr)\Bigr),
\]
sebab $(2\pi n)^{1/(2n)} = \exp\bigl(\frac{\ln(2\pi
n)}{2n}\bigr)$. Dengan membalikkannya, $(n!)^{-1/n} = \frac\eu n(1 +
O(\frac{\ln n}n))$, lalu menjumlahkannya atas $n \leq N$:
\[
\sum_{n\leq N}(n!)^{-1/n}
= \eu\ln N + O(1),
\qquad
\eu\sum_{n\leq N}\frac1n = \eu\ln N + O(1)
\]
(sebab deret galatnya $\sum\frac{\ln n}{n^2}$ konvergen). Sebuah
ketaksamaan Carleman berkonstanta $c$ akan memaksa $\eu\ln N
+ O(1) \leq c\,(\ln N + O(1))$, sehingga $c \geq \eu$ setelah
dibagi $\ln N$. Barisan pengoptimalnya sejajar:
sebab konstanta Hardy didekati oleh $n^{-1/p}$ (pertanyaan 21),
dan konstanta Carleman oleh $n^{-1}$ --- dalam setiap kasusnya barisan bertipe
harmonik yang duduk tepat di luar ruang yang dirata-ratakan.

\textbf{25.} Untuk $f = \mathbf 1_{\intcc01}$: $Hf(x) = 1$ pada
$\intoc01$ dan $Hf(x) = \frac1x$ untuk $x > 1$, sehingga
$\norm{Hf}_2^2 = 1 + \int_1^\infty x^{-2}\dd x = 2$ dan
nisbahnya $\sqrt2 \approx 1.414$, yakni sekitar $71\%$ batas
tajamnya. Sedangkan untuk $f_A$ (dengan $p = 2$): $F(x) = 2(\sqrt x - 1)$ pada
$\intcc1A$, sehingga pada rentang itu $Hf_A = 2x^{-1/2} - 2x^{-1}$
dan, untuk $x > A$, $Hf_A(x) = 2(\sqrt A - 1)/x$. Lalu dengan mengkuadratkan
dan mengintegralkannya,
\begin{align*}
\int_1^A\bigl(2x^{-1/2} - 2x^{-1}\bigr)^2\dd x
&= 4\ln A - 12 + \frac{16}{\sqrt A} - \frac4A,\\
\int_A^{\infty}\frac{4(\sqrt A - 1)^2}{x^2}\,\dd x
&= 4 - \frac{8}{\sqrt A} + \frac4A,
\end{align*}
sehingga $\norm{Hf_A}_2^2 = 4\ln A - 8 + 8A^{-1/2}$; lalu dibagi
$\norm{f_A}_2^2 = \ln A$ diperolehlah identitas yang dinyatakan itu. Di $A
= \eu^{10}$: $4 - \frac{8(1 - \eu^{-5})}{10} = 3.2054$, sehingga
nisbahnya $\sqrt{3.2054} \approx 1.790 < 2$. Sedangkan cacatnya
$4 - \norm{Hf_A}_2^2/\norm{f_A}_2^2 \sim 8/\ln A$ meluruh
hanya secara logaritmik: sehingga untuk mencapai nisbah $1.99$ orang memerlukan
$\ln A \approx 200$, yakni $A \approx 10^{87}$. Jadi konstanta
tajam adalah teorema, bukan percobaan.

\end{solution}
